\section{Spectral and coefficient transfers}\label{sec:reductions}
This appendix supplies the two transfers used to remove the auxiliary partition in Section~\ref{sec:estimate}.
All derivations in this appendix are relative: a coefficient is killed by the distinguished derivation and is not required to have a conjugate variable of its own.

\begin{lemma}\label{lem:lattice-transfer}
Let $P=P^*=P^2$ have trace $1/2$, and let $(e_i)$ be a finite partition of unity by nonzero projections with $t_i:=\tau(e_i)$.
On the word algebra $\mathcal N:=\operatorname{alg}^*(P,(e_i)_i)$ assume actual partial liberation
\[
 \gamma e_i=e_i\otimes1-1\otimes e_i,\qquad \gamma P=0,
 \qquad \eta:=\gamma^*(1\otimes1)\in L^2(N),
\]
where $N:=W^*(P,(e_i)_i)$.
Label the $e_i$ by distinct points $c_i$ of a square lattice of mesh $h>0$ and put
\begin{align*}
 U&:=\sum_ic_ie_i,\qquad
 d_c:=I_{\sum_{i\ne j}(c_i-c_j)^{-1}e_i\otimes e_j}\gamma,
 \qquad E_c:=(\id\otimes\tau)d_c,\\
 H_i&:=\sum_{j\ne i}\frac{t_j}{c_i-c_j},\qquad
 M_i:=\sup_{s\ge h}\frac{\sum_{|c_i-c_j|\le s}t_j}{s},\\
 B_c^2&:=\sum_it_i(M_i^2+|H_i|^2)
 +\sum_{i\ne j}\frac{\|e_i\eta e_j\|_2^2}{|c_i-c_j|^2}.
\end{align*}
There are the following transfers, with absolute constants.
\begin{enumerate}
\item Any such finite model admits an actual relative coordinate $V$ on disjoint circles of radius $h/8$ about $c_i$, with the same coefficient $P$, actual vacuum norm at most $cB_c$, and zero liberation mass on each circle square.
For every old word, its $d_c$ derivative is the conditional projection of its derivative in the enlarged model, interpreted by original-core graph approximation.
\item Rounding any actual relative coordinate $Z,P$ to spectral lattice cells gives models with $B_c\le c\|\xi\|_2$.
The word values converge in operator norm and their $d_c$ derivatives converge to the original derivatives in coarse Hilbert norm.
\end{enumerate}
Consequently a uniform slice bound for the circle models gives a uniform slice bound for all actual relative coordinates with coefficient $P$.
\end{lemma}
\begin{proof}
We first describe the reconstruction of a distinguished derivative from scalar and liberation data.
If $\mu$ is diffuse, has scalar vacuum $\xi_\mu$, and an actual liberation vacuum $\eta$ satisfies $\int|z-w|^{-2}\dd\nu_\eta<\infty$, then the relative derivative exists and
\begin{equation}\label{eq:reconstruct}
 \xi=\xi_\mu+S_{1/(\bar z-\bar w)}\eta,\qquad
 \|\xi\|_2^2=\|\xi_\mu\|_2^2+\int|z-w|^{-2}\dd\nu_\eta.
\end{equation}
To verify this, assign the normal phase to $Z^*$ and zero to the coefficients on formal words.
The identity $\gamma=I_{z-w}\delta$ and injectivity of the coarse inner gap make this prescription respect represented relations.
For $k_r(z,w):=(\bar z-\bar w)/(|z-w|^2+r^2)$, the marginal liberation adjoint formula gives, on every mixed word $p$,
\[
 \langle S_{\bar k_r}\eta+2\operatorname{row}(\bar k_r),p\rangle
 =\langle1\otimes1,I_{|z-w|^2/(|z-w|^2+r^2)}\delta p\rangle.
\]
The spectral term converges strongly by the inverse-gap condition.
The row terms are uniformly bounded by Lemma~\ref{lem:scalar}; testing marginal polynomials shows $2\operatorname{row}(\bar k_r)\rightharpoonup\xi_\mu$.
The coarse multiplier tends strongly to one.
This proves the actual vacuum identity on all mixed words.
The two terms in \eqref{eq:reconstruct} are orthogonal because the spectral term has zero diagonal component and $L^2(W^*(Z))$ is diagonal.

For the first transfer, freely adjoin independent standard semicircular operators $S_i$ to $N$.
In the normalized corner $e_i$, the operator $t_i^{-1/2}e_iS_ie_i$ is standard semicircular, and bounded Borel functions of it have conditional expectation onto $N$ equal to their normalized corner trace times $e_i$.
These facts follow from the semicircular moment and compression formulas \cite{NS06}; the conditional identity follows first for polynomials by the same moment calculation and then by bounded approximation.
Choose a Borel function of each diffuse corner variable to obtain $V_i$ uniform on the circle of radius $h/8$ about $c_i$ and put $V:=\sum_iV_i$.
Its component projections recover the $e_i$, so $N\subseteq M':=W^*(V,P)$.

On the old algebra let $\Lambda:=\operatorname{ad}-\gamma$, where $\operatorname{ad}(a):=a\otimes1-1\otimes a$.
It kills all $e_i$ and has vacuum $-\eta$.
Extend it by zero on the new free noises.
This preserves the actual vacuum: in every centered reduced word involving a new free factor, at least one trace leg at each cut retains a centered new letter, so its double trace vanishes.
Bounded functions of the compressed noises are obtained by bounded strong-star graph approximation, since the extended derivative kills both the $e_i$ and the noises.
Subtracting this extension from $\operatorname{ad}$ produces actual partial liberation at $V$, killing $P$, with vacuum $\eta$.
Also $e_i\eta e_i=0$: insertion of $e_i\otimes e_i$ into the old $\gamma$ is zero on generators and hence on all words, and its vacuum is $e_i\eta e_i$ because the two end corrections cancel.
The inverse-gap energy in the new model is at most
\[
 \frac{16}{9}\sum_{i\ne j}\frac{\|e_i\eta e_j\|_2^2}{|c_i-c_j|^2},
\]
since distinct circles are separated by at least $3|c_i-c_j|/4$.

Write $\mu_i'$ for uniform probability on the $i$th circle and $\mu':=\sum_it_i\mu_i'$.
Each circle has scalar vacuum $(z-c_i)/(h/8)^2$, of norm $8/h$.
This is checked on Laurent monomials by their circle moments.
The conjugate of the scalar vacuum for $\mu'$ on component $i$ is
\[
 t_i\overline{\xi_{\mu_i'}(z)}+2\sum_{j\ne i}t_j\int\frac1{z-w}\dd\mu_j'(w).
\]
The identity follows by separating the within-component and cross-component double integrals in the divided-difference pairing.
The cross kernel differs from $(c_i-c_j)^{-1}$ by at most $ch/|c_i-c_j|^2$.
Dyadic annuli and $t_i/h\le M_i$ bound the scalar vacuum squared norm by $c\sum_it_i(M_i^2+|H_i|^2)$.
Equation~\eqref{eq:reconstruct} supplies the actual relative derivative at $V$ with vacuum norm at most $cB_c$.

Smooth functions equal to one near one circle and zero near the others approximate the $e_i$ in operator norm and in the original derivative graph, using $C^1$ polynomial approximation.
Their derivative tensors have symbols $(1_i(z)-1_i(w))/(z-w)$.
On two distinct circles the product mean of $(z-w)^{-1}$ is exactly $(c_i-c_j)^{-1}$: expand in the centered displacements by a uniformly convergent geometric series and use the vanishing of every nonconstant circle moment.
The conditional expectation onto $N$ takes each bounded marginal function to its circle mean on each component.
Consequently its tensor square takes these graph derivatives to $d_ce_i$.
It also preserves the zero derivative of $P$.
Outer bimodularity recovers $d_cp$ for every old word, and projection of its right slice is contractive.
This proves the first transfer without presupposing domains beyond the original graph limits.

For the second transfer, partition the compact spectrum of $Z$ into half-open squares of side $h$, omit null cells and let $c_i$ be their centers.
Then $\|Z-U\|\le h$.
Liberation extends to bounded Borel marginal functions by bounded strong-star graph approximation, so its restriction to $N:=W^*(P,(e_i)_i)$ has vacuum $\eta=\E_N[Z^*,\xi]$.
Conditional expectation is bimodular over the cell algebra.
For $i\ne j$, old spectral distances are at most $3|c_i-c_j|$, so \eqref{eq:gap-budget} bounds the inverse-gap sum by $9\|\xi\|_2^2$.
For $x$ in cell $i$, $M_i\le4M_1\mu(x)$.
Comparison with the soft Cauchy transform at radius $h$ gives
\[
 \left|H_i-\int\frac{\overline{x-y}}{|x-y|^2+h^2}\dd\mu(y)\right|
 \le cM_1\mu(x).
\]
For cells within $6h$, both kernels are bounded by $1/h$ and their mass is at most $chM_1\mu(x)$.
On other cells their difference is at most $ch/|x-y|^2$; summing dyadic annuli gives the displayed estimate.
Lemma~\ref{lem:scalar} therefore yields $B_c\le c\|\xi\|_2$.
The generator tensors at $U,U^*,P$ are respectively
\[
 1\otimes1-\sum_ie_i\otimes e_i,\qquad
 \sum_{i\ne j}\frac{\bar c_i-\bar c_j}{c_i-c_j}e_i\otimes e_j,\qquad0.
\]
They converge almost everywhere off the diffuse spectral diagonal to the original generator tensors and have modulus at most one.
Dominated convergence, norm convergence of word values and the finite Leibniz expansion prove the asserted derivative convergence.
\end{proof}

\begin{lemma}\label{lem:projection-transfer}
Let $K_{\rm rel}$ be the supremum of $\|C\delta p\|_2/\|\xi\|_2$ over actual relative normal coordinates with finitely many bounded selfadjoint coefficients and core contractions $p$.
Let $K_{\rm proj}$ be the same supremum with one trace-half projection as coefficient.
In the extended nonnegative reals,
\[
 K_{\rm proj}\le K_{\rm rel}\le2K_{\rm proj}.
\]
\end{lemma}
\begin{proof}
Fix a relative model and one polynomial witness $p$.
Rescale its $m$ coefficients to contractions, set $k:=m+1$, $d:=2k$, $t:=1/d$, and freely adjoin standard semicirculars $s_2,\ldots,s_d$.
In the normalized $d$-by-$d$ matrix algebra put
\[
 W:=E_{11}\alpha Z+\sum_{r=2}^dE_{rr}(10r+s_r),\qquad
 B:=\sum_{r=2}^k(E_{1r}+E_{r1})(2+H_{r-1})^{1/2},
\]
where $0<\alpha\|Z\|\le1$ and $B$ is viewed in the upper $k$-by-$k$ corner.
Choose $c>\max(1,\|B\|)$ and form the trace-half projection
\[
 P_0:=\frac12
 \begin{pmatrix}I+B/c&(I-(B/c)^2)^{1/2}\\(I-(B/c)^2)^{1/2}&I-B/c\end{pmatrix}.
\]
The square roots are norm limits of coefficient polynomials, so their distinguished derivatives are zero in the corresponding graph.

We describe the derivative and its cost explicitly.
The entrywise lift of a base derivative $d_r$ sends a tensor $a\otimes b$ in entry $(i,j)$ to $E_{ir}a\otimes E_{rj}b$.
For normalized matrix trace its vacuum is $tE_{rr}\xi_r$, of norm $t^{3/2}\|\xi_r\|_2$.
Use $\delta/\alpha$ in slot one and the semicircular difference quotients in the other slots.
Their sum kills $P_0$ and takes $W$ to $\sum_rE_{rr}\otimes E_{rr}$.
Add the inner derivation with marginal kernel $(z-w)^{-1}$ on distinct summand spectra and zero on equal summands.
The spectra are separated by at least six, so this kernel has norm at most one, its vacuum costs at most two, and its derivative at $P_0$ has norm at most two.
The resulting $L$ satisfies $LW=1\otimes1$.

Freely adjoin one more standard semicircular $S$ and put $Y:=P_0+\varepsilon S$, where $2\varepsilon<1/4$.
Extend $L$ by zero on $S$ and add $I_{-LP_0/\varepsilon}\partial_S$.
It kills $Y$ and leaves the values at $W,W^*$ unchanged.
Its added vacuum has norm $\|LP_0\|_2/\varepsilon$: on base tensors the map giving this vacuum is the free isometry $a^*\otimes b^*\mapsto a^*Sb^*$.
The reduced-word calculation proves the actual vacuum identities for these free extensions.
Now $P:=1_{(1/2,\infty)}(Y)$ is a norm limit of polynomials in $Y$ and has zero derivative.
The spectral gap and continuous functional calculus give $\|P-P_0\|\to0$ as $\varepsilon\downarrow0$.
For small $\varepsilon$ these projections are equivalent, hence $\tau(P)=1/2$.
Restricting to $W^*(W,P)$ and projecting the vacuum supplies an actual relative model with
\begin{equation}\label{eq:encoding-vacuum}
 \|\xi'\|_2\le t^{3/2}\|\xi\|_2/\alpha+t^{3/2}\sqrt{d-1}+2+2/\varepsilon.
\end{equation}
The forced normal phase makes the restricted generator tensors intrinsic to this algebra.

Let $e:=E_{11}$ and $e_r:=E_{rr}$.
The corner coefficients
\[
 H_l':=(2c)^2ePe_{l+1}Pe-2e
\]
converge in norm to $E_{11}H_l$, since the $(1,l+1)$ entry of $P_0$ is $(2+H_l)^{1/2}/(2c)$.
All $e_r$ are smooth spectral cuts of $\Rea W$, constant on neighborhoods of the separated components.
Polynomial $C^1$ approximation places them in the original relative graph with derivative symbol $(e_r(z)-e_r(w))/(z-w)$.
Insert $e\otimes e$ in both inner legs and denote the resulting derivative by $\theta$.
It kills the $e_r$, $P$ and the $H_l'$, and its value on the first corner is the original phase after multiplication by $\alpha$.

If $a'$ is the slice norm of this projection model, the product rule gives for $x=exe$ in this graph algebra
\[
 (\id\otimes\tau')I_{e\otimes e}\delta' x
 =(E'(xe)-xE'(e))e,\qquad E':=(\id\otimes\tau')\delta'.
\]
Its norm is at most $2a'\|x\|$.
Under the normalized corner trace $\tau_e=\tau'/t$, the left side has norm $t^{3/2}\|(\id\otimes\tau_e)\theta x\|_{2,\tau_e}$.
Evaluate the fixed word $p$ at $eWe/\alpha,H_1',\ldots,H_m'$ with unit $e$.
As $\varepsilon\downarrow0$, its norm tends to $\|p\|\le1$, and its scaled derivative slice tends to $E_{11}C\delta p$ in normalized corner norm.
These comparisons are independent of $\alpha$ after multiplication of the derivative by $\alpha$.
Thus, for any small $\eta>0$, first choose $\varepsilon$ so that
\[
 2\alpha a'(1+\eta)\ge t^{3/2}(\|C\delta p\|_2-\eta).
\]
Next send $\alpha$ to zero in \eqref{eq:encoding-vacuum} and then $\eta$ to zero.
The ratio $a'/\|\xi'\|_2$ approaches from below at least $\|C\delta p\|_2/(2\|\xi\|_2)$.
If $a'$ is infinite the same conclusion for the suprema is immediate.
Taking the supremum over the original finite witnesses proves the upper comparison; the lower comparison is inclusion of classes.
\end{proof}
